
\begin{titlepage}
    % Strona tytułowa
    \vbox to\textheight{\hyphenpenalty=10000
    \begin{center}
	\begin{tabular}{p{107mm} p{9cm}}
	    \begin{minipage}{9cm}
	      \begin{center}
	      Politechnika Warszawska \\
	      Wydział Elektroniki i~Technik Informacyjnych \\
	      Instytut Informatyki
	      \end{center}
	    \end{minipage}
	    &
	    \begin{minipage}{8cm}
	    \begin{flushleft}
	     \footnotesize
	      Rok akademicki 2013/2014
	    \vspace*{2.75\baselineskip}
	    \end{flushleft}
	    \end{minipage} \\
	\end{tabular}
        \vspace*{2.5\baselineskip}
	\par\vspace{\smallskipamount}
	\begin{center}
		\includegraphics[angle=0,scale=1]{img/logo_pw.png}
	\end{center}
	\vspace*{2\baselineskip}{\LARGE Praca dyplomowa inżynierska\par}
	\vspace{3\baselineskip}{\LARGE\strut Krzysztof Opasiak\par}
	\vspace*{2\baselineskip}{\huge\bfseries Rozproszony monitoring systemów komputerowych\par}

	\vspace*{7\baselineskip}
	\hfill\mbox{}\par\vspace*{\baselineskip}\noindent
	\begin{tabular}[b]{@{}p{3cm}@{\ }l@{}}
	    {\large\hfill } & {\large }
	\end{tabular}
	\hfill
	\begin{tabular}[b]{@{}l@{}}
	Opiekun pracy: \\[\smallskipamount]
	{\large dr inż. Piotr Gawkowski}
	\end{tabular}\par
	\vspace*{4\baselineskip}
    \begin{tabular}{p{\textwidth}}
    \begin{flushleft}
	\begin{minipage}{7cm}
	Ocena \dotfill
	\par\vspace{1.6\baselineskip}
	\dotfill
	\par\noindent
	\centerline{\footnotesize Podpis Przewodniczącego} \par
	\centerline{\footnotesize Komisji Egzaminu Dyplomowego}\par
	\end{minipage}
    \end{flushleft}
    \end{tabular}
    \end{center}}

    % Życiorys
    \newpage\thispagestyle{empty}
    \begin{tabular}{p{4cm} p{12cm}}
    \begin{minipage}{5cm}
    
    \includegraphics[height=5cm,width=4cm]{img/foto.jpg}
    \end{minipage}
    &
    \begin{minipage}{12cm}
    \begin{flushleft}
    \par\noindent\vspace{1\baselineskip}
    \begin{tabular}[h]{l l}
    {\normalsize\it Kierunek:} & {\normalsize \hspace{0.8cm} Informatyka}
    \end{tabular}\par\noindent\vspace{1\baselineskip}
    \begin{tabular}[h]{l l}
    {\normalsize\it Specjalność:}& {\normalsize \hspace{0.4cm} Inżynieria Systemów Informatycznych}
    \end{tabular}
    \par\noindent\vspace{1\baselineskip}
    \begin{tabular}[h]{l l}
    {\normalsize\it Data urodzenia:} & {\normalsize \hspace{4.6cm} 1990.12.28}
    \end{tabular}
    \par\noindent\vspace{1\baselineskip}
    \begin{tabular}[h]{l l}
    {\normalsize\it Data rozpoczęcia studiów:} & {\normalsize \hspace{2.8cm} 2010.10.01}
    \end{tabular}
    \par\noindent\vspace{1\baselineskip}
    \end{flushleft}
    \end{minipage}
    \end{tabular}
    \vspace*{1\baselineskip}
    \begin{center}
	{\large\bfseries Życiorys}\par\bigskip
    \end{center}

    \indent 
    Urodziłem się 28 grudnia 1990 roku w Koninie. Uczęszczałem
    kolejno do Szkoły Podstawowej numer 8 im. Powstańców
    Wielkopolskich w Koninie, a następnie Gimnazjum Towarzystwa
    Salezjańskiego w Koninie.

    \par
    \indent 
    W latach 2006-2010 uczęszczałem do Technikum w Zespole
    Szkół im. Mikołaja Kopernika w Koninie. W trakcie nauki w tej
    szkole dwukrotnie przyznano mi stypendium Prezesa Rady Ministrów
    za bardzo dobre wyniki w nauce oraz wzorowe zachowanie.  W roku
    2010 ukończyłem z wyróżnieniem szkołę średnią, a następnie zdałem
    maturę oraz egzamin zawodowy uzyskując tytuł Technik
    Teleinformatyk.

    \par
    \indent
    W październiku 2010 roku rozpocząłem studia stacjonarne pierwszego stopnia na Wydziale Elektroniki
    i Technik Informacyjnych na kierunku Informatyka.

    \vspace{2\baselineskip}
    \hfill\parbox{15em}{{\small\dotfill}\\[-.3ex]
    \centerline{\footnotesize podpis studenta}}\par
    \vspace{2\baselineskip}
    \begin{center}
 	{\large\bfseries Egzamin dyplomowy} \par\bigskip\bigskip
    \end{center}
    \par\noindent\vspace{1.5\baselineskip}
    Złożył egzamin dyplomowy w dn. \dotfill 20\_\_r
    \par\noindent\vspace{1.5\baselineskip}
    Z wynikiem \dotfill
    \par\noindent\vspace{1.5\baselineskip}
    Ogólny wynik studiów \dotfill
    \par\noindent\vspace{1.5\baselineskip}
    Dodatkowe wnioski i uwagi Komisji \dotfill
    \par\noindent\vspace{1.5\baselineskip}
    \dotfill

    % Streszczenie
    \newpage\thispagestyle{empty}
    \vspace*{2\baselineskip}
    \begin{center}
	{\large\bfseries Streszczenie}\par\bigskip
    \end{center}

    {\itshape Przedmiotem tej pracy jest system monitorowania
      rozproszone uwzględniający wymagania dotyczące monitorowania
      klienta mobilnego.

      \indent Na początku pracy wykonany został przegląd dostępnych na
      rynku systemów monitorujących. Następnie przedstawiono
      charakterystykę klienta mobilnego i~statycznego oraz
      zdefiniowano wymagania stawiane przed projektowanym
      systemem. Opisano również architekturę systemu Icinga, a także
      omówiono dostępne do niego dodatki.

      \indent Kolejna część zawiera zawiera opis proponowanej przez
      autora architektury rozwiązania. W szczególny sposób opisany
      został zaprojektowany w ramach tej pracy protokół komunikacyjny.

      \indent Przedstawiono również opis wykonanej implementacji
      dodatku umożliwiającego przekazywanie w~bezpieczny sposób danych
      z~urządzenia mobilnego do systemu Icinga.

      \indent W~końcowej części tej pracy znajduje się opis wykonanej
      konfiguracji testowej zaprojektowanego systemu. Zawarto tam
      również krótkie sprawozdanie oraz spostrzeżenia z~okresu
      testowego użytkowania systemu.}\vspace*{1\baselineskip}

    \noindent{\bf Słowa kluczowe}: {\itshape monitorowanie,
      monitorowanie rozproszone, Icinga, protokół komunikacyjny, bezpieczeństwo komunikacji.}
    \par
    \vspace{4\baselineskip}
    \begin{center}
	{\large\bfseries Abstract}\par\bigskip
    \end{center}
    \noindent{\bf Title}: {\itshape Distributed monitoring of computer
      systems.}\par
    \vspace*{1\baselineskip} {\itshape The subject of this thesis is
      distributed monitoring system which takes into account
      requirements of mobile devices.

      \indent In the beginning of this thesis review of few available
      on market monitoring systems has been presented. Then
      characteristic of mobile and stationary clients has been
      introduced and after that, requirements of proposed
      system. Detailed architecture of Icinga and few available
      addons has also been presented.

      \indent Next part contains description of proposed
      solution. Definition of communication protocol developed in this
      thesis has also been pointed out there.

      \indent Implementation of proposed addon has been
      described. This addon allow to transport data from mobile device
      to Icinga in a secure manner.

      \indent End of this thesis contains of description of test
      configuration of proposed system and short report with author remarks
      of use of the created system. } \vspace*{1\baselineskip}

    \noindent{\bf Key words}: {\itshape monitoring, distributed
      monitoring, Icinga, communiaction protocol, communication security.}

\end{titlepage}

% ex: set tabstop=4 shiftwidth=4 softtabstop=4 noexpandtab fileformat=unix filetype=tex spelllang=pl,en spell:
